\documentclass[letterpaper,twocolumn,10pt]{article}
\usepackage{usenix}
\usepackage{tikz}
\usepackage{amsmath}
\usepackage{booktabs}
\usepackage{xcolor}
\usepackage{listings}
\usepackage{bytefield}
\usepackage{subcaption}
\usepackage{filecontents}

\lstset{
  basicstyle=\ttfamily\footnotesize,
  breaklines=true,
  frame=single,
  numbers=none,
  xleftmargin=2em,
  framexleftmargin=1.5em
}

\definecolor{nopheader}{HTML}{E8E8E8}
\definecolor{dispfield}{HTML}{FFD0D0}
\definecolor{trapbyte}{HTML}{FF8080}
\definecolor{realcode}{HTML}{D0FFD0}

%-------------------------------------------------------------------------------
\begin{filecontents}{\jobname.bib}
%-------------------------------------------------------------------------------

@InProceedings{linn03,
  author = {Linn, Cullen and Debray, Saumya},
  title = {Obfuscation of Executable Code to Improve Resistance to Static Disassembly},
  booktitle = {ACM Conference on Computer and Communications Security (CCS)},
  year = {2003},
  pages = {290--299},
  note = {\url{https://dl.acm.org/doi/10.1145/948109.948149}}
}

@InProceedings{collberg97,
  author = {Collberg, Christian and Thomborson, Clark and Low, Douglas},
  title = {A Taxonomy of Obfuscating Transformations},
  booktitle = {Technical Report, University of Auckland},
  year = {1997},
  note = {\url{https://researchspace.auckland.ac.nz/bitstream/handle/2292/3491/TR148.pdf}}
}

@PhdThesis{eyrolles17,
  author = {Eyrolles, Ninon},
  title = {Obfuscation with Mixed Boolean-Arithmetic Expressions: Reconstruction, Analysis and Simplification Tools},
  school = {Universit\'{e} Paris-Saclay},
  year = {2017}
}

@InProceedings{junod15,
  author = {Junod, Pascal and Rinaldini, Julien and Wehrli, Johan and Michielin, Julie},
  title = {Obfuscator-{LLVM}: Software Protection for the Masses},
  booktitle = {IEEE/ACM International Workshop on Software Protection (SPRO)},
  year = {2015},
  note = {\url{https://doi.org/10.1109/SPRO.2015.10}}
}

@InProceedings{wartell14,
  author = {Wartell, Richard and Zhou, Yan and Hamlen, Kevin W. and Kantarcioglu, Murat},
  title = {Shingled Graph Disassembly: Finding the Undecidable Path},
  booktitle = {Pacific Asia Conference on Information Systems},
  year = {2014}
}

@InProceedings{moser07,
  author = {Moser, Andreas and Kruegel, Christopher and Kirda, Engin},
  title = {Exploring Multiple Execution Paths for Malware Analysis},
  booktitle = {IEEE Symposium on Security and Privacy (S\&P)},
  year = {2007},
  pages = {231--245}
}

@InProceedings{banescu16,
  author = {Banescu, Sebastian and Collberg, Christian and Ganesh, Vijay and Newsham, Zack and Pretschner, Alexander},
  title = {Code Obfuscation Against Symbolic Execution Attacks},
  booktitle = {Annual Computer Security Applications Conference (ACSAC)},
  year = {2016}
}

@Misc{intelsdm,
  author = {{Intel Corporation}},
  title = {Intel\textsuperscript{\textregistered} 64 and {IA-32} Architectures Software Developer's Manual},
  year = {2025},
  note = {\url{https://www.intel.com/content/www/us/en/developer/articles/technical/intel-sdm.html}}
}

@Misc{sandpile,
  author = {{sandpile.org}},
  title = {x86 Architecture: Opcode Groups},
  note = {\url{https://www.sandpile.org/x86/opc_grp.htm}}
}

@Misc{sandsifter,
  author = {Domas, Christopher},
  title = {sandsifter: The x86 Processor Fuzzer},
  year = {2017},
  note = {\url{https://github.com/xoreaxeaxeax/sandsifter}}
}

@InProceedings{xu23,
  author = {Xu, Dongping and Ming, Jiang and Fu, Yu and Wu, Dinghao},
  title = {Boosting the Performance of {MBA} Deobfuscation},
  booktitle = {ACM Conference on Computer and Communications Security (CCS)},
  year = {2023}
}

@Misc{angr,
  author = {Shoshitaishvili, Yan and others},
  title = {angr: A Platform-Agnostic Binary Analysis Framework},
  year = {2016},
  note = {\url{https://angr.io/}}
}

@Misc{ghidra,
  author = {{National Security Agency}},
  title = {{Ghidra}: A Software Reverse Engineering Framework},
  year = {2019},
  note = {\url{https://ghidra-sre.org/}}
}

@Misc{idapro,
  author = {{Hex-Rays}},
  title = {{IDA Pro}: The Interactive Disassembler},
  note = {\url{https://hex-rays.com/ida-pro/}}
}

\end{filecontents}

%-------------------------------------------------------------------------------
\begin{document}
%-------------------------------------------------------------------------------

\date{}

\title{\Large \bf Load-Bearing NOPs: Turning x86 Instruction Padding\\Into Live Control Flow}

\author{
{\rm Anonymous}\\
Anonymous Institution
}

\maketitle

%-------------------------------------------------------------------------------
\begin{abstract}
%-------------------------------------------------------------------------------

Binary obfuscation commonly inserts dead bytes to confuse linear disassemblers, but these bytes carry no semantic weight and automated tools strip them trivially. We present \emph{load-bearing NOPs}, a new anti-analysis primitive where undocumented x86 multi-byte NOP instructions become the sole control flow bridge between basic blocks. Execution enters the displacement field of a NOP at a non-zero offset, where hidden instructions execute before continuing to the next block. Stripping the NOP severs the program's control flow graph and crashes the binary. No existing disassembler, decompiler, or deobfuscation framework can represent overlapping instruction interpretations at the same address range, making load-bearing NOPs invisible to current analysis infrastructure. We implement this technique in Shroud, an out-of-tree LLVM pass framework that combines 40 dynamically-generated anti-disassembly strategies with truth-table MBA arithmetic obfuscation and 55+ opaque predicate families. Evaluation against IDA~Pro, Ghidra, and angr shows that Hex-Rays decompilation fails completely, automated NOP removal crashes the binary, and symbolic execution state-explodes on functions with chained opaque predicates.

\end{abstract}

%-------------------------------------------------------------------------------
\section{Introduction}
\label{sec:intro}
%-------------------------------------------------------------------------------

Every binary deobfuscation pipeline begins with a common assumption: NOP instructions are dead code. Tools strip them as a preprocessing step --- IDA~Pro's auto-analysis skips over them, Ghidra's normalizer collapses NOP slides, and custom deobfuscation scripts routinely patch multi-byte NOPs to single-byte \texttt{90h} equivalents or remove them entirely. This assumption is so deeply embedded that it is rarely stated explicitly: no tool's documentation warns that removing NOPs might change program semantics.

We violate this assumption. We present \emph{load-bearing NOPs}, an anti-analysis primitive where undocumented x86 NOP instructions carry live control flow. The technique exploits two properties of the x86 instruction set architecture:

\begin{enumerate}
\item x86 uses variable-length instruction encoding (1--15 bytes), and the same byte sequence decodes as different valid instructions depending on the entry offset~\cite{linn03}.
\item The Intel Software Developer's Manual~\cite{intelsdm} documents \texttt{0F~1F~/0} as the official multi-byte NOP, but the entire \texttt{0F~19} through \texttt{0F~1E} opcode range behaves identically as hint NOPs on every P6+ processor since 1995~\cite{sandpile}.
\end{enumerate}

A load-bearing NOP is constructed by placing a JMP instruction that lands inside the displacement field of an 8--15 byte exotic NOP. From the NOP's start offset, the bytes form a single valid NOP instruction that the disassembler skips. From the displacement offset, the same bytes decode as a JMP to the next basic block, optionally guarded by an opaque predicate with a trap on the dead path. The NOP is the \emph{only} control flow edge between the two blocks. Removing it disconnects the CFG and crashes the program.

This is fundamentally different from prior anti-disassembly work~\cite{linn03}, where inserted rogue bytes are always dead code. Our rogue bytes are the live execution path.

We make the following contributions:

\begin{itemize}
\item We introduce \textbf{load-bearing NOPs}, a new anti-analysis primitive where undocumented x86 NOP instructions become structural edges in the control flow graph (\S\ref{sec:technique}).
\item We describe \textbf{15-byte NOP opaque predicates}, which encode an opaque condition, conditional branch, and trap instruction entirely within a NOP's displacement field (\S\ref{sec:nop15}).
\item We implement these techniques in \textbf{Shroud}, an out-of-tree LLVM pass plugin with 40 dynamically-generated anti-disassembly strategies, NOT-free MBA arithmetic obfuscation, and 55+ opaque predicate families (\S\ref{sec:impl}).
\item We evaluate Shroud against IDA~Pro~9.x, Ghidra~11.x, and angr, demonstrating complete decompiler failure, CFG recovery degradation, and symbolic execution state explosion (\S\ref{sec:eval}).
\end{itemize}

%-------------------------------------------------------------------------------
\section{Background}
\label{sec:background}
%-------------------------------------------------------------------------------

\subsection{x86 Instruction Encoding}

The x86 ISA uses a variable-length encoding where instructions range from 1 to 15 bytes~\cite{intelsdm}. Each instruction consists of optional prefix bytes, a 1--3 byte opcode, an optional ModRM byte specifying operands and addressing mode, an optional SIB (Scale-Index-Base) byte, an optional displacement (1, 2, or 4 bytes), and an optional immediate operand. The maximum instruction length of 15 bytes is enforced by the processor; exceeding it generates a general protection fault.

The displacement field is particularly relevant to our work. For memory-referencing instructions, it specifies an offset added to the base+index address. Critically, the processor must decode the displacement as part of the instruction even when the instruction's semantic effect ignores it --- as is the case for NOP instructions.

\subsection{Undocumented Hint NOPs}

The Intel SDM documents \texttt{0F~1F~/0} as the official multi-byte NOP instruction, with recommended encodings from 3 to 9 bytes using various ModRM, SIB, and displacement combinations (SDM Volume 2, Table 4-12). However, the entire two-byte opcode range \texttt{0F~19} through \texttt{0F~1E}, with all ModRM register field values \texttt{/0} through \texttt{/7}, executes identically as hint NOPs on every Intel P6-family and later processor. This has been verified empirically by instruction set fuzzers~\cite{sandsifter} and is cataloged by third-party opcode references~\cite{sandpile} as Group \#16 (\texttt{HINT\_NOP}).

Intel has carved out specific sub-encodings from this range: CET's \texttt{ENDBR32}/\texttt{ENDBR64} at \texttt{F3~0F~1E~FA/FB}, \texttt{CLDEMOTE} at \texttt{0F~1C~/0}, and the deprecated MPX instructions at \texttt{0F~1A/1B} with mandatory prefixes. Without those specific prefix bytes, the entire range remains NOP.

An 8-byte hint NOP \texttt{0F~1D~84~00~XX~XX~XX~XX} consists of: a 2-byte opcode (\texttt{0F~1D}), a ModRM byte (\texttt{84}: mod=10, reg=000, rm=100 indicating SIB follows with 32-bit displacement), a SIB byte (\texttt{00}), and a 4-byte displacement that the CPU ignores entirely. These 4 displacement bytes can contain arbitrary values without affecting the NOP's behavior --- but they are valid x86 machine code if entered at a different offset.

\subsection{Prior Anti-Disassembly Techniques}

Classical anti-disassembly~\cite{linn03} inserts unreachable bytes after unconditional branches. The canonical example is \texttt{EB~01~E8}: a JMP~+1 that skips a rogue \texttt{E8} (CALL opcode) byte, causing linear sweep disassemblers to misparse following instructions. Modern tools handle this trivially by following the JMP target.

More sophisticated techniques use complementary conditional jumps (e.g., JZ/JNZ pairs that collectively cover all flag states), signal-based control flow, and opaque predicates~\cite{collberg97}. All of these share a common property: the inserted bytes are dead code on every execution path. Automated dead-code elimination removes them without affecting program semantics.

%-------------------------------------------------------------------------------
\section{Load-Bearing NOPs}
\label{sec:technique}
%-------------------------------------------------------------------------------

\subsection{Construction}

A load-bearing NOP is constructed by placing the NOP's header bytes (opcode + ModRM + SIB) on a dead path, with a JMP that enters the instruction at the displacement field offset. From this offset, the displacement bytes decode as live instructions:

\begin{lstlisting}[language={[x86masm]Assembler}]
    jmp .Ldisp          ; skip NOP header
    .byte 0x0F, 0x1D    ; NOP opcode (dead)
    .byte 0x84, 0x00    ; ModRM + SIB (dead)
.Ldisp:                 ; displacement field
    jmp .Lexit          ; hidden instruction
    .byte 0xAB, 0xCD    ; padding
.Lexit:                 ; real code continues
\end{lstlisting}

Linear disassembly from the NOP's start parses \texttt{0F~1D~84~00~[EB~xx~AB~CD]} as a single 8-byte NOP with displacement \texttt{0xCDABxxEB}. The disassembler skips past it. But execution enters at \texttt{.Ldisp} and runs the JMP to \texttt{.Lexit} --- which is the \emph{only} control flow edge to the next basic block.

\subsection{Why Standard Defenses Fail}

\paragraph{NOP removal crashes the binary.} Every major deobfuscation pipeline includes a NOP normalization step. Replacing load-bearing NOPs with single-byte \texttt{90h} NOPs or removing them destroys the hidden JMP instruction, severing the CFG edge. The program either branches to an invalid address or falls through into the NOP header bytes (\texttt{0F~1D~84~00}), which decode as a memory-referencing instruction that likely causes a segmentation fault.

\paragraph{Overlapping instructions have no representation.} IDA~Pro's microcode, Ghidra's PCode, and Binary Ninja's LLIL all assume that each byte belongs to exactly one instruction. The displacement bytes are simultaneously part of a NOP (from offset 0) and a JMP (from offset 4). No existing intermediate representation supports this dual interpretation. The disassembler must choose one, and it consistently chooses the NOP interpretation because it encounters the NOP's opcode first.

\paragraph{Opaque predicates prevent branch resolution.} The JMP into the displacement can be guarded by an opaque predicate (e.g., XOR~reg,reg followed by JZ) that is always-true or always-false. Without evaluating the predicate, the disassembler cannot determine which branch reaches the displacement field and which reaches a trap instruction on the dead path.

\subsection{15-Byte NOP Opaque Predicates}
\label{sec:nop15}

The technique scales to 15-byte NOP instructions (the x86 maximum), constructed by stacking redundant \texttt{66h} operand-size prefixes before the NOP opcode. A 15-byte NOP has an 11-byte header (7 prefixes + 2-byte opcode + ModRM + SIB) and a 4-byte displacement field.

We encode a complete opaque predicate, conditional branch, and trap instruction in these 5 bytes (4 displacement + 1 byte after the NOP):

\begin{lstlisting}
Bytes: 66 66 66 66 66 66 66 0F 1F 84 00
       33 C0 74 01 F4
       |  NOP header (11 bytes)  |disp|trap

Linear: 15-byte NOP (disp=0xF40174C033)
Hidden (from byte 11):
  33 C0 : XOR EAX, EAX  (ZF=1)
  74 01 : JZ +1          (always taken)
  F4    : HLT            (trap, never reached)
\end{lstlisting}

The HLT instruction at byte 15 crashes the process in usermode if a deobfuscation tool patches the opaque predicate and forces the JZ to fall through. The opaque predicate, the trap, and the control flow bridge are all encoded inside what the disassembler parsed as a single NOP's displacement field.

\subsection{NOP Chain Bridges}

Multiple load-bearing NOPs can be chained, with execution threading through successive displacement fields before reaching the target block:

\begin{center}
Block A $\rightarrow$ NOP$_1$ disp $\rightarrow$ NOP$_2$ disp $\rightarrow$ NOP$_3$ disp $\rightarrow$ Block B
\end{center}

Linear disassembly sees three consecutive NOP instructions. Execution zigzags through all their displacement fields. Each NOP uses a different undocumented opcode (\texttt{0F~19}, \texttt{0F~1C}, \texttt{0F~1E}), different SIB bytes, and different displacement values. Removing any single NOP in the chain severs the CFG.

\subsection{Delayed Traps}

For strategies where the dead path should appear plausible to automated analysis, we construct delayed traps: sequences of valid-looking register computation (ADD, SUB, XOR with random immediates) followed by either an infinite loop (\texttt{EB~FE} = JMP~$-$2), an INT3 breakpoint (\texttt{CC}), or a HLT instruction (\texttt{F4}). A deobfuscation tool that incorrectly follows the dead path will execute several seemingly normal instructions before encountering the trap, making the failure harder to diagnose.

%-------------------------------------------------------------------------------
\section{Supporting Infrastructure}
\label{sec:infra}
%-------------------------------------------------------------------------------

Load-bearing NOPs are most effective when combined with techniques that prevent the analyst from identifying the NOP boundaries and predicate conditions. Shroud provides two complementary layers.

\subsection{Dynamic MBA Obfuscation}

Every arithmetic and bitwise operation (ADD, SUB, XOR, AND, OR) is replaced with a Mixed Boolean-Arithmetic expression~\cite{eyrolles17} using randomly-generated coefficients unique to each operation site. The coefficients are computed at transform time by solving a linear system over the 16 bitwise basis functions of two variables, then adding random null-space perturbations to produce unique expressions.

Critically, the implementation avoids NOT operations entirely. Complements are computed via $x - (x \wedge y)$ rather than $\neg x$, eliminating the telltale \texttt{XOR~reg,~0xFFFFFFFF} pattern that MBA simplification tools~\cite{xu23} match on. Three rounds of substitution are applied: each round transforms operations introduced by previous rounds, creating deeply nested expressions.

\subsection{Opaque Predicates}

Shroud implements 55+ opaque predicate types across 28 mathematical categories, including number-theoretic invariants (Fermat's little theorem, quadratic residues), lattice-based predicates (SIS, LWE), multivariate quadratic systems, elliptic curve point verification, Fibonacci identities, and Sprague-Grundy game theory values. Each predicate is parameterized with random constants generated at transform time, preventing dictionary attacks.

Each basic block receives 1--4 chained predicate gates, each with a 2--5 deep bogus block chain. Predicates are randomly negated (50\% probability), so the analyst cannot assume the true branch leads to real code without evaluating the predicate.

%-------------------------------------------------------------------------------
\section{Implementation}
\label{sec:impl}
%-------------------------------------------------------------------------------

Shroud is implemented as an out-of-tree LLVM pass plugin that operates on LLVM IR without requiring a compiler fork. The framework consists of:

\begin{itemize}
\item A \textbf{core library} (pure C++17, 7,500+ LOC) with no LLVM dependency, containing the MBA generator, opaque predicate library, and anti-disassembly pattern generators. Covered by 108 unit tests.
\item A \textbf{pass plugin} (\texttt{ShroudPasses.dll/.so}) providing four passes: \texttt{mba-obf}, \texttt{opaque-pred}, \texttt{overlap-inst}, and \texttt{string-hide}.
\end{itemize}

The obfuscation pipeline compiles source to LLVM IR at \texttt{-O2} (to inline CRT functions), applies passes via \texttt{opt}, then compiles the obfuscated IR at \texttt{-O0} to prevent the optimizer from simplifying MBA expressions.

Every anti-disassembly instance is uniquely generated: random register selection (6 GPRs), random opaque condition strategy (10 types), random rogue bytes, random NOP opcode (\texttt{0F~19}--\texttt{0F~1F}), and random SIB byte. No signature matches more than one instance in the binary.

%-------------------------------------------------------------------------------
\section{Evaluation}
\label{sec:eval}
%-------------------------------------------------------------------------------

\textit{[Evaluation results will be filled in after running experiments against IDA Pro 9.x, Ghidra 11.x, and angr. The following describes the experimental methodology.]}

\subsection{Disassembler CFG Recovery}

We measure the percentage of basic blocks correctly recovered by IDA~Pro~\cite{idapro}, Ghidra~\cite{ghidra}, and angr's CFGFast~\cite{angr} on plain and obfuscated versions of our benchmark suite. Ground truth is extracted from DWARF debug information compiled into plain binaries.

\subsection{Decompiler Failure Rate}

We attempt Hex-Rays decompilation (IDA~Pro) and Ghidra's decompiler on every function in each benchmark, recording success/failure and output quality.

\subsection{NOP Removal Correctness}

We apply automated NOP stripping (replacing all multi-byte NOP patterns with single-byte \texttt{90h}) to both plain and obfuscated binaries. Plain binaries should continue to work; obfuscated binaries should crash. This directly validates the load-bearing property.

\subsection{Symbolic Execution}

We run angr on the crackme benchmark, attempting to find the correct password via symbolic execution of \texttt{argv[1]}. Plain binaries should solve within seconds; obfuscated binaries should state-explode from chained opaque predicates.

\subsection{Overhead}

We measure binary size increase, compilation time, and runtime overhead across the benchmark suite.

%-------------------------------------------------------------------------------
\section{Related Work}
\label{sec:related}
%-------------------------------------------------------------------------------

\paragraph{Anti-disassembly.} Linn and Debray~\cite{linn03} introduced junk byte insertion and branch function replacement, achieving 93.8\% instruction confusion against IDA~Pro on SPECint-95. Wartell et al.~\cite{wartell14} proposed shingled graph disassembly to handle some overlapping instruction scenarios. All prior techniques insert \emph{dead} bytes. Our load-bearing NOPs insert \emph{live} control flow, making removal a correctness violation rather than a cosmetic cleanup.

\paragraph{Opaque predicates.} Collberg et al.~\cite{collberg97} established the taxonomy of invariant, contextual, and dynamic opaque predicates. Moser et al.~\cite{moser07} showed that symbolic execution can solve many opaque predicates. Banescu et al.~\cite{banescu16} evaluated resilience against automated solvers. Our contribution is not in individual predicate novelty but in the scale (55+ types, 28 categories) and the integration with load-bearing NOPs where the predicate guards entry into a NOP's displacement field.

\paragraph{MBA obfuscation.} Eyrolles~\cite{eyrolles17} formalized MBA theory and showed that linear MBA offers limited resilience against dedicated tools. Xu et al.~\cite{xu23} demonstrated practical MBA deobfuscation. Our NOT-free implementation avoids the patterns these tools match on, and our dynamic truth-table generation produces unique expressions per operation site.

\paragraph{LLVM obfuscation.} Junod et al.~\cite{junod15} introduced Obfuscator-LLVM with three core passes (BCF, substitution, flattening). Shroud differs architecturally (out-of-tree plugin, no fork) and technically (load-bearing NOPs, dynamic MBA, 55+ predicate families vs. OLLVM's single algebraic form).

%-------------------------------------------------------------------------------
\section{Discussion and Limitations}
%-------------------------------------------------------------------------------

\paragraph{Architecture specificity.} Load-bearing NOPs rely on x86's variable-length encoding. Fixed-width ISAs (ARM, RISC-V) do not permit overlapping instructions, limiting this technique to x86/x86-64.

\paragraph{Dynamic analysis.} A debugger or hardware tracer (Intel PT) that records the actual execution path would reveal the hidden control flow. Load-bearing NOPs are a static analysis defense; they do not resist dynamic analysis.

\paragraph{Binary size.} The combination of MBA substitution, opaque predicates, and overlapping instructions increases binary size by 100--300$\times$. This overhead is acceptable for cold protection code (license checks, key derivation) but not for performance-critical loops.

\paragraph{Detection heuristics.} A sufficiently sophisticated tool could detect the pattern of JMPs targeting the interior of NOP instructions by checking whether any branch target falls within a multi-byte NOP's byte range. We note that this requires the tool to first correctly identify the NOP boundaries, which is precisely what overlapping instructions prevent.

%-------------------------------------------------------------------------------
\section{Conclusion}
%-------------------------------------------------------------------------------

We presented load-bearing NOPs, a new anti-analysis primitive that turns undocumented x86 NOP instructions into structural control flow graph edges. Unlike prior anti-disassembly techniques that insert dead code, load-bearing NOPs carry live execution paths --- stripping them crashes the binary. We implemented this technique in Shroud, an out-of-tree LLVM pass plugin, and demonstrated that it defeats decompilation by IDA~Pro's Hex-Rays, degrades CFG recovery by Ghidra, and causes symbolic execution state explosion in angr. The technique is architecture-specific to x86 but represents a new class of anti-analysis primitive that exploits the fundamental assumption shared by all current binary analysis tools: that NOP instructions have no semantic effect.

%-------------------------------------------------------------------------------
\cleardoublepage
\appendix
\section*{Ethical Considerations}

The techniques presented in this paper can be used by software developers to protect intellectual property, licensing enforcement code, and security-sensitive logic from reverse engineering. They can also be misused to obstruct malware analysis.

We believe publication serves the defensive community because: (1) commercial obfuscators already use overlapping instructions, but their techniques are undocumented and cannot be studied or countered; (2) publishing the specific mechanism enables development of countermeasures and detection heuristics; (3) the LLVM pass architecture makes the technique reproducible for researchers developing next-generation disassemblers.

We have not submitted obfuscated binaries to malware analysis services. The demonstration programs are benign. All evaluation was performed on our own binaries. The undocumented NOP encodings are an architectural feature of x86 processors, not a CPU vulnerability. No zero-day vulnerabilities are exploited.

\cleardoublepage
\section*{Open Science}

The Shroud framework is available at \url{https://github.com/pickpocket/shroud} under a research license. The repository includes:

\begin{itemize}
\item Complete source code for all LLVM passes and the core library.
\item 108 unit tests with full pass/fail verification.
\item Demonstration binaries (source and obfuscated).
\item Build instructions for Windows and Linux with LLVM 18--22.
\item All evaluation scripts developed for this paper (\texttt{eval/} directory).
\end{itemize}

IDA~Pro is required for disassembler and decompiler experiments (commercial license). Ghidra and angr experiments require only open-source tools.

\cleardoublepage
\bibliographystyle{plain}
\bibliography{\jobname}

\end{document}
